\chapter{LX-009 --- Localize Shell Help, Notification, and Directional Surfaces}\label{record-356}

\noindent{\footnotesize\color{BookMuted}\path{docs/method/graveyard/legacy-backlog/LX-009-localize-shell-help-notification-and-directional-surfaces.md}\par}\smallskip

Legend: \href{https://github.com/flyingrobots/bijou/blob/cbd70612140e7a450d764b3c7be1c6af878b99bc/docs/method/graveyard/legends/LX-localization-and-bidirectionality.md}{LX --- Localization and Bidirectionality}

\section{Idea}\label{record-356-idea}

LX-008 proved the shell i18n seam in real product surfaces:

\begin{itemize}
\tightlist
\item
  frame-owned shell copy can come from catalogs
\item
  DOGFOOD can load localized shell and app copy
\item
  shell direction metadata can drive at least one visible mirrored behavior
\end{itemize}

The next slice should finish the most obvious remaining gaps:

\begin{itemize}
\tightlist
\item
  help overlay binding descriptions and grouping copy
\item
  notification-center history/body strings and filters
\item
  broader logical-direction behavior in shell-owned drawers and related chrome
\end{itemize}

\section{Why}\label{record-356-why}

LX-008 intentionally stopped once the i18n seam was real.

That was the correct stopping point, but it leaves visible English islands and only partial direction-aware shell behavior. LX-009 should tighten those gaps before we claim the shell is broadly localized.

\section{Status}\label{record-356-status}

Backlog spawned by the retrospective for LX-008.

